\documentclass{article}

% Language setting
\usepackage[english]{babel}

% Set page size and margins
\usepackage[letterpaper,top=2cm,bottom=2cm,left=3cm,right=3cm,marginparwidth=1.75cm]{geometry}

% Useful packages
\usepackage{amsmath}
\usepackage{amssymb}
\usepackage{graphicx}
\usepackage[colorlinks=true, allcolors=black]{hyperref}
\usepackage{titlesec}
\usepackage{xparse}

\NewDocumentCommand{\qcq}{m o o m}{%
	\ensuremath{\forall #1 \IfValueT{#2}{, #2} \IfValueT{#3}{, #3}  \in  \mathbb{#4}}%
}
\NewDocumentCommand{\ext}{m o o m}{%
	\ensuremath{\exists #1 \IfValueT{#2}{, #2} \IfValueT{#3}{, #3}  \in  \mathbb{#4}}%
}

\title{\textbf{Arithmétique dans $\mathbb{Z}$}}
\author{\textbf{Yassin Akkab}}

\newcommand{\R}{\mathbb{R}}
\newcommand{\N}{\mathbb{N}}
\newcommand{\Z}{\mathbb{Z}}
\newcommand{\C}{\mathbb{C}}
\newcommand{\Q}{\mathbb{Q}}
\newcommand{\D}{\mathbb{D}}

\begin{document}
	
	\maketitle
	\tableofcontents
	\newpage
	
	\section{Divisibilité et division euclidienne}
	
	\subsection{Divisibilité dans $\mathbb{Z}$}
	Soient $a$ et $b$ deux entiers relatifs. On dit que $a$ \textbf{divise} $b$ (ou que $b$ est un \textbf{multiple} de $a$), et on note $a \mid b$, s'il existe un entier relatif $k$ tel que :
	\[ b = k \cdot a \]
	Propriétés :
	\begin{itemize}
		\item Si $a \mid b$ et $b \mid c$, alors $a \mid c$.
		\item Si $a \mid b$ et $a \mid c$, alors pour tous entiers $\alpha$ and $\beta$, $a \mid (\alpha b + \beta c)$.
		\item Si $a \mid b$ et $b \ne 0$, alors $|a| \le |b|$.
	\end{itemize}
	
	\subsection{Division euclidienne}
	\textbf{Théorème} : Pour tout entier relatif $a$ et tout entier naturel non nul $b$, il existe un unique couple d'entiers relatifs $(q, r)$ tel que :
	\[ a = bq + r \quad \text{avec} \quad 0 \le r < b \]
	$q$ est le quotient et $r$ est le reste de la division euclidienne de $a$ par $b$.
	
	\section{PGCD et PPCM}
	
	\subsection{Plus Grand Commun Diviseur (PGCD)}
	Soient $a$ et $b$ deux entiers non tous deux nuls. Le plus grand diviseur commun positif de $a$ et $b$ est noté $\text{pgcd}(a, b)$ ou $a \wedge b$.
	\begin{itemize}
		\item Deux entiers $a$ et $b$ sont dits \textbf{premiers entre eux} si $\text{pgcd}(a, b) = 1$.
		\item \textbf{Algorithme d'Euclide} : Si $a = bq + r$, alors $\text{pgcd}(a, b) = \text{pgcd}(b, r)$. On répète les divisions euclidiennes jusqu'à obtenir un reste nul. Le dernier reste non nul est le PGCD.
	\end{itemize}
	
	\subsection{Plus Petit Commun Multiple (PPCM)}
	Soient $a$ et $b$ deux entiers non nuls. Le plus petit multiple commun strictement positif de $a$ et $b$ est noté $\text{ppcm}(a, b)$ ou $a \vee b$.
	\textbf{Relation fondamentale} :
	\[ (a \wedge b) \cdot (a \vee b) = |ab| \]
	
	\section{Théorèmes de Bézout et de Gauss}
	
	\subsection{Théorème de Bézout}
	Soient $a$ et $b$ deux entiers relatifs non nuls.
	\[ \text{pgcd}(a, b) = d \iff \ext{u}[, v][]{Z}, \quad au + bv = d \]
	En particulier, $a$ et $b$ sont premiers entre eux si et seulement s'il existe deux entiers relatifs $u$ et $v$ tels que :
	\[ au + bv = 1 \]
	
	\subsection{Théorème de Gauss}
	Soient $a$, $b$ et $c$ trois entiers relatifs non nuls.
	\[ \text{Si } a \mid bc \quad \text{et} \quad \text{pgcd}(a, b) = 1, \quad \text{alors } a \mid c \]
	
	\subsection{Équations diophantiennes $ax + by = c$}
	Soit l'équation $(E) : ax + by = c$ d'inconnues $(x, y) \in \mathbb{Z}^2$.
	\begin{itemize}
		\item L'équation $(E)$ admet des solutions dans $\mathbb{Z}^2$ si et seulement si $c$ est un multiple de $\text{pgcd}(a, b)$.
		\item Si $(x_0, y_0)$ est une solution particulière de $(E)$ et $d = \text{pgcd}(a, b)$, avec $a = a'd$ et $b = b'd$, la solution générale est :
		\[ x = x_0 + kb', \quad y = y_0 - ka' \quad (\text{où } k \in \mathbb{Z}) \]
	\end{itemize}
	
	\section{Congruences dans $\mathbb{Z}$}
	
	\subsection{Définition et compatibilité}
	Soit $n$ un entier naturel supérieur ou égal à 2. Deux entiers $a$ et $b$ sont dits \textbf{congrus modulo $n$}, et on note $a \equiv b \pmod n$, si $a - b$ est un multiple de $n$.
	La relation de congruence est une relation d'équivalence compatible avec l'addition et la multiplication :
	Si $a \equiv b \pmod n$ et $c \equiv d \pmod n$, alors :
	\begin{itemize}
		\item $a + c \equiv b + d \pmod n$
		\item $a \cdot c \equiv b \cdot d \pmod n$
		\item $\qcq{k}{N}, \quad a^k \equiv b^k \pmod n$
	\end{itemize}
	
	\section{Nombres premiers et Fermat}
	
	\subsection{Théorème fondamental de l'arithmétique}
	Tout entier naturel $n \ge 2$ s'écrit de manière unique sous la forme d'un produit de facteurs premiers :
	\[ n = p_1^{\alpha_1} p_2^{\alpha_2} \dots p_k^{\alpha_k} \]
	où $p_1 < p_2 < \dots < p_k$ sont des nombres premiers et $\alpha_i \in \mathbb{N}^*$.
	
	\subsection{Petit théorème de Fermat}
	Soit $p$ un nombre premier et $a$ un entier relatif.
	\begin{itemize}
		\item De manière générale : $a^p \equiv a \pmod p$.
		\item Si $p$ ne divise pas $a$ (c'est-à-dire $\text{pgcd}(a, p) = 1$) :
		\[ a^{p-1} \equiv 1 \pmod p \]
	\end{itemize}
	
\end{document}
